% from ega4-18-fr.tex
\oldpage[IV-4]{131}

  \item[a')] La condition $a)$ est satisfaite pour toutes les $A$-algèbres $B$ de la forme $A[T]/F.A[T]$, où $F\in A[T]$ est un polynôme unitaire.
  \item[b)] Soient $S=\Spec(A)$, $S_0=\Spec(k)$. Pour tout morphisme étale $g:S'\to S$, si l'on pose $S'_0=S'\otimes_A k$, toute $S_0$-section $u_0$ de $S'_0$ est la restriction d'une $S$-section $u$ de $S'$.
  \item[c)] Pour tout morphisme $f:X\to S$, séparé et localement de type fini, et tout point $x\in X$ tel que $f(x)$ soit égal au point fermé $s$ de $S$ et que $f$ soit quasi-fini au point $x$ (Err$_{\mathrm{III}}$, 20), $X$ est somme de deux préschémas $X'$, $X''$ tels que $X'=\Spec(\sh{O}_{X,x})$ et que $f|X':X'\to S$ soit un morphisme fini.
  \item[c')] Pour tout morphisme localement de type fini $f:X\to S$, et tout point $x\in X$ tel que $f$ soit quasi-fini au point $x$, et que $f(x)$ soit égal au point fermé $s$ de $S$, $\sh{O}_{X,x}$ est une algèbre finie sur $\sh{O}_{S,s}=A$.
  \item[c'')] Pour tout morphisme localement de présentation finie $f:X\to S$, et tout point $x\in X$ tel que $f$ soit quasi-fini au point $x$, et que $f(x)$ soit égal au point fermé $s$ de $S$, $\sh{O}_{X,x}$ est une algèbre finie et de présentation finie sur $\sh{O}_{S,s}=A$.
\end{enumerate}
\end{theorem}

Notons d'abord que la condition $c')$ (resp. $c'')$ est équivalente à la même condition où on suppose de plus $f$ \emph{séparé}, la question étant locale sur $X$. De même, la condition $b)$ est équivalente à la même condition où l'on suppose de plus $g$ \emph{séparé}: en effet, il suffit d'appliquer cette dernière à la restriction de $g$ à un voisinage ouvert affine du point $u_0(S_0)$ dans $S'$. Bornons-nous donc désormais au cas où, dans $b)$, $c')$ et $c'')$, les morphismes donnés sont \emph{séparés}.

Le fait que $a)$ implique $b)$ et que $b)$ implique $a')$ est un cas particulier de (18.5.4). Montrons en outre que $a')$ implique $a)$. Il s'agit de prouver que si $e_0$ est un idempotent de $C=B\otimes_A k$, il existe un idempotent $e\in B$ dont $e_0$ est l'image canonique. Si $b$ est un élément de $B$ dont l'image dans $C$ est $e_0$, la sous-$A$-algèbre $A[b]=B'$ de $B$ est finie, et l'image canonique $C'$ de $B'\otimes_A k$ dans $C$ contient $e_0$. Or, $B'\otimes_A k$ est une $k$-algèbre finie, donc composée directe de $k$-algèbres locales finies, et par suite $e_0$ est l'image dans $C'$ d'un idempotent $e'_0$ de $B'\otimes_A k$. On est ainsi ramené au cas où $B$ est monogène, et par suite isomorphe à une $A$-algèbre quotient d'une algèbre de la forme $A[T]/F.A[T]$, où $F$ est un polynôme unitaire. Or, en vertu de $a')$, $A[T]/F.A[T]$ est composée directe d'anneaux locaux, donc il en est de même de toutes ses algèbres quotients, ce qui achève de prouver l'existence de l'idempotent $e$.

Il est immédiat que $c)$ entraîne $a)$, comme on le voit en raisonnant par récurrence sur le nombre d'idéaux maximaux de l'anneau semi-local $B$. Pour voir que $a)$ implique $c)$, on peut, en vertu de (13.1.4), se borner au cas où $f$ est un morphisme affine et quasi-fini. Alors, par application du « Main theorem » (8.12.8), $f$ peut s'écrire comme un morphisme composé
\[
  X\xrightarrow{j}Y\xrightarrow{g}S,
\]
où $g$ est un morphisme fini et $j$ une immersion ouverte. Comme $Y=\Spec(B)$, où $B$ est une $A$-algèbre finie, il résulte de $a)$ que $B$ est composée directe d'anneaux locaux, qui sont évidemment des $A$-algèbres finies, et $\sh{O}_{X,x}$ s'identifie à l'un de ces anneaux locaux puisque $f(x)=s$; en outre, tout ouvert de $Y$ contenant $x$ contient nécessairement $\Spec(\sh{O}_{X,x})$.

Il est trivial que $c)$ implique $c')$, en vertu de la remarque du début. Prouvons que $c')$ implique $c'')$. Supposons en effet $c')$ vérifiée, et prouvons que sous les conditions

